\chapter{A deep dive into \mytool}\label{chap:approach}

% --------------------- Section 1 ---------------------

\section{Introduction}\label{sec:introduction}
In Chapter~\ref{chap:running_example}, we presented a practical example showcasing the usage of \mytool{}.
This chapter dives deeper into its design principles and implementation details, elucidating the underlying mechanisms that enable stable specialization through metaprogramming.
The actual implementation is in Rust, but as stated earlier, the concepts can be generalized to other programming languages with similar capabilities.
First, we begin by outlining the core design goals that guided the development of the tool.

\paragraph{Design Goals}
The primary design goals for \mytool{} were centered around achieving effective specialization in stable code while adhering to the principles of safety, usability and performance.
These goals shaped the tool's architecture, macro design and code generation techniques.
The key objectives included:
\begin{itemize}
	\item \textbf{Stable Compatibility}: Ensure that the tool operates entirely within the confines of stable Rust, avoiding any reliance on nightly features or unstable APIs.
	\item \textbf{Type Safety}: Maintain Rust's strong type safety guarantees, ensuring that specialized implementations are correctly typed and validated at compile time.
	\item \textbf{Ease of Use}: Provide a straightforward and intuitive interface for developers to define specialized implementations without excessive boilerplate or complexity.
	\item \textbf{Performance Efficiency}: Generate specialized code that is optimized for performance, minimizing overhead and ensuring that the benefits of specialization are realized in practice.
\end{itemize}

\paragraph{Overview}
As already seen, \mytool{} leverages the macro system to enable specialization.
The three primary components needed by the approach are:
\begin{itemize}
	\item \textbf{\nameref{sec:conditions_definition}}: A macro that annotates specialized implementations, specifying the conditions under which they should be used and generates the necessary boilerplate.
	\item \textbf{\nameref{sec:method_wrapping}}: A second macro that wraps method calls, directing them to the appropriate specialized implementation.
	\item \textbf{\nameref{sec:build_script}}: A way to analyze the code at compile time and collect the necessary information to facilitate specialization.
\end{itemize}

% --------------------- Section 2 ---------------------

\section{Conditions Definition}\label{sec:conditions_definition}
The first step to achieve specialization through metaprogramming is to define various specialized implementations that should be applied when certain conditions are met.
In \mytool{}, this is accomplished using the \inlinerust{#[when(predicate)]} attribute macro, which takes as argument the predicate, which is the condition that must be satisfied for the specialized implementation to be selected.
Please note that if we wanted a simpler drop-in replacement for the unstable specialization feature, we could edit the tool to avoid passing the predicate altogether and just use normal rust syntax for implementations, but we still would need to mark the implementation with some kind of macro to generate the necessary boilerplate.

\paragraph{Predicates}
In our implementation, the macro supports these predicates:
\begin{itemize}
	\item \textbf{Type Equality}: Check if a generic type parameter is of a specific type, e.g.\@ \mbox{\texttt{T = i32}}.
	\item \textbf{Trait Implementation}: Check if a generic type parameter implements a specific trait or multiple traits, e.g.\@ \mbox{\texttt{T:\@ Display}} or \mbox{\texttt{T:\@ Debug + Clone}}.
	\item \textbf{Logical And}: Check if all the provided predicates are satisfied, e.g.\@ \mbox{\inlinerust{all(p1, p2, ...)}}. % TODO: fix formatting
	\item \textbf{Logical Or}: Check if at least one of the provided predicates is satisfied, e.g.\@ \mbox{\inlinerust{any(p1, p2, ...)}}.
	\item \textbf{Logical Not}: Check if the provided predicate is not satisfied, e.g.\@ \mbox{\inlinerust{not(predicate)}}.
\end{itemize}
These predicates can be combined to form complex conditions, allowing for flexible and powerful specialization rules.

\paragraph{Wildcards}
In addition to the specific predicates, \mytool{} supports a wildcard condition that matches any type.
For instance, the condition \mbox{\texttt{T = Vec<\_>}} matches any vector type, regardless of its element type.
Another useful wildcard condition is \mbox{\texttt{T = \&\_}}, which matches any reference type.

\paragraph{Default Implementation}
In addition to specialized implementations, \mytool{} supports a default implementation that is used when no other conditions are met.
These default implementations are picked up by default by the \nameref{sec:build_script} and don't need to be explicitly marked with any macro.
Optionally, if we wanted to limit the implementations available for specialization, we could use another macro to explicitly mark the default implementations.

\paragraph{Normalization}
To ensure that the conditions are consistent and can be easily compared, \mytool{} normalizes the conditions.
This normalization process involves transforming the conditions into a canonical form, which allows for easier comparison and resolution of overlapping conditions.
The normal form chosen is a disjunctive normal form (DNF~\cite{DNFPost21}), where each condition is represented as a conjunction of predicates.
This means that each condition is a logical OR of one or more predicates, and each predicate is a logical AND of one or more atomic conditions.
The normalization process also involves applying De Morgan's laws~\cite{DeMorgan} to push negations down to the atomic conditions, ensuring that the conditions are expressed in a consistent manner.
This representation allows for efficient checking of whether a condition is satisfied, because we can easily check each disjunctive clause separately.
For example, the condition \texttt{not(all(\mbox{T:\@ Display}, any(\mbox{T = i32}, \mbox{T = f64})))} is normalized to \texttt{any(not(\mbox{T:\@ Display}), all(not(\mbox{T = i32}), not(\mbox{T = f64})))}.

\paragraph{Macro expansion}
Whenever possible, the conditions are evaluated at compile time and the generic parameters are replaced or constrained accordingly.
For example: if we have a type equality condition, the generic parameter is replaced with the specific type; if we have a trait implementation condition, the generic parameter is constrained to implement the specified trait(s).
Given that all the conditions are normalized, the macro expansion can be greatly optimized: the \texttt{when} macro expands each disjunctive condition into a separate specialized implementation, allowing for efficient resolution of the conditions at compile time.
This is the reason why we previously defined this approach \textit{monomorphic specialization}.

% TODO: fix positioning
\begin{figure}[htbp!]
	\centering
	\showrust*[0.98\linewidth]{conditions-definition-pre.rs}
	\captionsetup{type=listing,skip=2pt}
	\caption{Example of the conditions definition using the macro}
	\label{fig:conditions-definition-pre}
\end{figure}

\begin{figure}[htbp!]
	\centering
	\showrust*[0.98\linewidth]{conditions-definition-post.rs}
	\captionsetup{type=listing,skip=2pt}
	\caption{Example of the expansion of the conditions definition defined in Listing \ref{fig:conditions-definition-pre}}
	\label{fig:conditions-definition-post}
\end{figure}

% --------------------- Section 3 ---------------------

\section{Method Wrapping}\label{sec:method_wrapping}
To enable the selection of the appropriate specialized implementation at method call sites, we need a second macro, which wraps method calls and directs them to the correct implementation.
In \mytool{}, this is accomplished using the \inlinerust{spec!} macro, which takes as first arguments the method call itself.

\paragraph{Type Annotations}
If the language supported reflection or if we had some way to inspect the type of the arguments at compile time, this would have been enough.
However, Rust does not support such features, so we need to provide additional information to the macro.
This additional information includes the types of object on which the method is called and the types of the arguments passed to the method, if any.
\begin{itemize}
	\item \textbf{Type of the Object}: The type of the object on which the method is called, e.g.\@ if we have \inlinerust{x.fn(arg1, arg2)}, we need to provide the type of \texttt{x}, e.g.\@ \mbox{\texttt{MyType}}.
	\item \textbf{Types of the Arguments}: A list of the types of the arguments passed to the method, if the method takes no arguments, it can be empty or omitted, e.g.\@ if we have \inlinerust{x.fn(arg1, arg2)}, we need to provide the types of \texttt{arg1} and \texttt{arg2}, e.g.\@ \mbox{\texttt{[Arg1Type, Arg2Type]}}.
\end{itemize}
With this information, the tool can solve almost all of the specialization at compile time, generating a direct call to the appropriate specialized implementation.

\paragraph{Annotations}
The \textit{type annotations} defined above are the first two of what we call \textit{annotations}, which are additional pieces of information provided to the macro to help it resolve the specialization.
In Rust, just like there is no way to inspect the type of a variable at compile time, there is also no way to inspect the traits implemented by a type or to inspect type aliases.
To overcome these limitations, \mytool{} allows the user to provide additional annotations to the method wrapping macro.
These additional annotations include:
\begin{itemize}
	\item \textbf{Implemented Traits}: A list of traits implemented by each relevant type, e.g.\@ \mbox{\texttt{i32:\@ Display}} or \mbox{\texttt{i32:\@ Debug + Clone}}.
	\item \textbf{Type Aliases}: A list of type aliases for each relevant type, e.g.\@ \mbox{\texttt{MyInt = i32}}.
\end{itemize}
These annotations can be provided as additional arguments to the \inlinerust{spec!} macro (without any particular order), allowing the tool to accurately resolve the specialization even in the presence of type aliases or when traits are involved in the conditions.
Considering all of this, the full syntax for the \inlinerust{spec!} macro is the following: \inlinerust{spec!{ method_call; var_type; args_types; annotations... }}, where \texttt{method\_call} is the method call, \texttt{var\_type} is the type of the object on which the method is called, \texttt{args\_types} is a list of the types of the arguments passed to the method (if any), and \texttt{annotations} are any additional annotations provided by the user.
For languages that support class inheritance, additional annotations could be added to reflect parent classes, e.g.\@ \mbox{\texttt{T1 <: T2}}.

\begin{figure}[htbp!]
	\centering
	\showrust*[0.98\linewidth]{method-wrapping-pre.rs}
	\captionsetup{type=listing,skip=2pt}
	\caption{Example of method wrapping, using implementations defined in Listing \ref{fig:conditions-definition-pre}}
	\label{fig:method-wrapping-pre}
\end{figure}

\begin{figure}[htbp!]
	\centering
	\showrust*[0.98\linewidth]{method-wrapping-post.rs}
	\captionsetup{type=listing,skip=2pt}
	\caption{Example of the expansion of method wrapping from Listing \ref{fig:method-wrapping-pre}}
	\label{fig:method-wrapping-post}
\end{figure}

% --------------------- Section 4 ---------------------

\section{Build Script}\label{sec:build_script}
% collect traits, default impls, handle order by collecting all first
\lipsum[1-4]

\paragraph{Alternatives}\label{subsec:alternatives}
% using #[specializable] and #[spec_default]
% drawback: not handling order
\lipsum[1-4]

% --------------------- Section 5 ---------------------

\section{Lifetimes}\label{sec:lifetimes}
% initially not supported, now supported in both #[when] and spec!
\lipsum[1-4]

% --------------------- Section 6 ---------------------

\section{Overlap}\label{sec:overlap}
% conditions in dnf, handling overlapping conditions
\lipsum[1-4]

% --------------------- Section 7 ---------------------

\section{Limitations}\label{sec:limitations}
\lipsum[1-4]